% =====================================================================
%  STANDALONE REVIEW DOCUMENT  --  NOT PART OF THE THESIS
%
%  Proposed edits introducing a general definition of
%  "ramification bounded by r" at a prime divisor (Z, xi).
%
%  Compile FROM THE PROJECT ROOT so that \input{includes} resolves:
%     latexmk -pdf -outdir=ramification_bound_general \
%         ramification_bound_general/ramification_bound_general.tex
% =====================================================================
\documentclass[11pt]{report}
\usepackage{geometry}\geometry{margin=1.05in}
\usepackage[T1]{fontenc}
\usepackage{ifthen}

\setlength{\parindent}{1em}
\setlength{\parskip}{0.5em}

\input{includes}
\input{MacrosAndFigures}

\usepackage{framed}

% ---------------------------------------------------------------------
% Review conventions
% ---------------------------------------------------------------------
\definecolor{newmath}{rgb}{0.00,0.20,0.75}
\definecolor{oldgray}{rgb}{0.35,0.35,0.35}
\newcommand{\new}[1]{\textcolor{newmath}{#1}}
\newcommand{\loc}[1]{\par\noindent{\small\textsf{#1}}\par}

\newenvironment{current}%
  {\par\medskip\noindent{\small\textbf{Currently:}}\par\nopagebreak
   \begingroup\color{oldgray}\small\leftskip=1.2em\relax}%
  {\par\endgroup\medskip}

\newenvironment{proposed}%
  {\par\medskip\noindent{\small\textbf{Proposed:}}\par\nopagebreak
   \begin{leftbar}}%
  {\end{leftbar}\medskip}

% ---------------------------------------------------------------------
% Phantom targets for labels that live elsewhere in main.tex.
% Numbers copied from main.aux so the cross-references read correctly.
% ---------------------------------------------------------------------
\makeatletter
\@namedef{r@definition:tamely_ramified_dvrs}{{35}{20}{}{definition.35}{}}
\@namedef{r@definition:tamely_ramified_dvrs@cref}{{[definition][35][]35}{[1][20][]20}{}{}{}}
\@namedef{r@definition:tamely_ramified_in_codim_1}{{38}{22}{}{definition.38}{}}
\@namedef{r@definition:tamely_ramified_in_codim_1@cref}{{[definition][38][]38}{[1][22][]22}{}{}{}}
\@namedef{r@cor:abelian_equivilance_fundamental_torsors}{{14}{11}{}{cor.14}{}}
\@namedef{r@cor:abelian_equivilance_fundamental_torsors@cref}{{[cor][14][]14}{[1][11][]11}{}{}{}}
\@namedef{r@section:ramification_of_discrete_valuations}{{2.1}{20}{}{section.2.1}{}}
\@namedef{r@section:ramification_of_discrete_valuations@cref}{{[section][1][2]2.1}{[1][20][]20}{}{}{}}
\@namedef{r@section:ramification_of_torsors}{{2.2}{22}{}{section.2.2}{}}
\@namedef{r@section:ramification_of_torsors@cref}{{[section][2][2]2.2}{[1][22][]22}{}{}{}}
\@namedef{r@section:cyclic_extension_theories}{{2.4}{30}{}{section.2.4}{}}
\@namedef{r@section:cyclic_extension_theories@cref}{{[section][4][2]2.4}{[1][30][]30}{}{}{}}
\makeatother

\title{A General Definition of \emph{Ramification Bounded by $r$}\\[0.3em]
       {\large Proposed edits to \S2.1--\S2.2 --- review copy}}
\author{}
\date{\today}

\begin{document}
\maketitle

\begin{abstract}
\noindent
This is a standalone review document. It is not part of the thesis and is not
included by \texttt{main.tex}. It records four proposed edits whose purpose is
to make the phrase ``ramification bounded by $r$ at $(Z,\xi)$'' --- used in the
statements and proofs of the stability lemmas of
\Cref{section:ramification_of_torsors}, and applied at the generic point of an
exceptional divisor --- into a defined notion. Material that is
\new{new or changed} is set in blue. Theorem and definition numbers inside this
document are local to it and bear no relation to the numbering in
\texttt{main.pdf}.
\end{abstract}

% =====================================================================
\chapter*{The gap}
\addcontentsline{toc}{chapter}{The gap}
% =====================================================================

\Cref{section:ramification_of_torsors} opens in the right generality: a locally
Noetherian scheme $X$, a dense open $U\subset X$, and a prime divisor $(Z,\xi)$
of $X\setminus U$ whose local ring $\Oo_{X,\xi}$ is a discrete valuation ring.
In that generality it defines only \emph{unramified} and \emph{tamely ramified}
(\Cref{definition:tamely_ramified_in_codim_1}). It then says ``now back to our
original assumptions'' and descends to a closed point $P$ of the curve $C$,
where it defines \emph{ramification bounded by $r$} and the character
dictionary $\rho_P:G_{\widehat K_P}\to G$.

Two things are therefore used but never introduced in the generality in which
they are used:

\begin{enumerate}
    \item the phrase \emph{$\cP$ has ramification bounded by $r$ at $(Z,\xi)$},
    which is the hypothesis and the conclusion of the contracted-product
    lemma and of the three stability lemmas that follow it;
    \item the character reformulation
    $\rho_\xi\bigl(G^{\,r}_{\widehat K_\xi}\bigr)=1$ at a general $\xi$, which
    is what the proofs of those lemmas actually manipulate --- each of them
    begins ``let $A=\widehat{\Oo_{X,\xi}}$, $K=\Frac(A)$, and let
    $\rho_i:G_K\to G$ be the associated characters''.
\end{enumerate}

This is not a cosmetic gap. The lemmas are invoked at $\eta_\ndls$, the generic
point of the exceptional divisor $E_\ndls\subset X_\ndls$ --- exactly the case
the closed-point definition does not cover.

\section*{The design question, and the choice made here}

The upper-numbering filtration is currently constructed in
\Cref{section:ramification_of_discrete_valuations} under the standing
hypothesis \emph{complete, equal characteristic $p$, perfect residue field}.
At $\eta_\ndls$ the residue field is $k(E_\ndls)=k(\bP^{d-1})$, a function
field, hence \emph{imperfect} in characteristic $p$. So the general definition
cannot simply cite that setup as it stands. Three ways out:

\begin{enumerate}
    \item Require $\kappa(\xi)$ perfect. Cleanest to state --- and excludes the
    only case we need.
    \item Develop ramification theory for imperfect residue fields
    (Abbes--Saito). Correct, and out of all proportion to the need.
    \item \new{Weaken the standing hypothesis from ``perfect residue field'' to
    ``separable residue extension''.}
\end{enumerate}

Option 3 is what is proposed below, for three reasons. It is Serre's own
standing hypothesis in \cite[Chapter~IV]{serreLF}; every construction in
\Cref{section:ramification_of_discrete_valuations} goes through under it
verbatim, since the only place perfectness is used is to produce a monogenic
generator of $B$ over $A$; and it is already what the proofs verify by hand ---
the unramifiedness argument in the Artin--Schreier case establishes precisely
that each field factor has ``ramification index $1$ and separable residue
extension''.

The cost is that residual separability becomes a hypothesis one must carry and
check. That is honest: it \emph{is} a hypothesis at an imperfect residue field,
and pretending otherwise is what produced the gap in the first place.

% =====================================================================
\chapter*{Edit A --- weaken the standing hypothesis}
\addcontentsline{toc}{chapter}{Edit A --- standing hypothesis}
% =====================================================================
\loc{File \texttt{sections/Preliminaries/LocalRamificationPreliminaries.tex}, replacing lines 45--55.}

\begin{current}
For the remainder of this section, assume that $A$ is a discrete valuation ring
of equal characteristic $p>0$, complete with respect to its
maximal-ideal-adic topology, with uniformizer $\pi$ and \textbf{perfect residue
field} $\kappa$. Thus $\operatorname{char}(A)=\operatorname{char}(\kappa)=p$.
The Cohen structure theorem provides a coefficient field $k\subset A$ mapping
isomorphically onto $\kappa$. Consequently, $A\cong k[[\pi]]\cong k[[t]]$, and
its fraction field is $K\cong k((\pi))$. Let $L/K$ be a finite Galois extension
with Galois group $G$. The integral closure $B$ of $A$ in $L$ is itself a
complete discrete valuation ring. Since the residue field $\kappa$ is perfect,
the extension of residue fields is separable, and there exists a uniformizer
$\Pi\in B$ such that $B=A[\Pi]$.
\end{current}

\begin{proposed}
For the remainder of this section, assume that $A$ is a discrete valuation
ring, complete with respect to its maximal-ideal-adic topology (cf.\
\stackstag{00IE} and \stackstag{0324}), with uniformizer $\pi$, residue field
$\kappa$ and fraction field $K$. Let $L/K$ be a finite Galois extension with
Galois group $G$, and let $B$ be the integral closure of $A$ in $L$; it is
again a complete discrete valuation ring. \new{We assume throughout that}
\begin{equation}\label{eq:residually_separable}
    \new{\text{the residue extension } \kappa_B/\kappa \text{ is separable.}}
\end{equation}
\new{This is the standing hypothesis of \cite[Chapter~IV]{serreLF}, and it is
what makes the constructions below available: under
\eqref{eq:residually_separable} there is an element $x\in B$ with $B=A[x]$
\cite[Chapter~III, \S6, Proposition~12]{serreLF} (one may take $x$ to be a
uniformizer $\Pi$ when $L/K$ is totally ramified). Condition
\eqref{eq:residually_separable} holds automatically whenever $\kappa$ is
perfect, and in particular when $\kappa$ is finite or algebraically closed; but
it is strictly weaker, and the extra generality is what allows us to speak of
ramification at the generic point of an exceptional divisor, whose residue
field is a function field and hence imperfect in characteristic $p$.}

\new{In the equal-characteristic situation with $\kappa$ perfect --- the
setting of \Cref{section:cyclic_extension_theories} --- one has more:} if
$\operatorname{char}(A)=\operatorname{char}(\kappa)=p>0$, the Cohen structure
theorem (\stackstag{032A}; see also \stackstag{0C0S}) provides a coefficient
field $k\subset A$ mapping isomorphically onto $\kappa$, so that
$A\cong k[[\pi]]\cong k[[t]]$ and $K\cong k((\pi))$.
\end{proposed}

\paragraph{Knock-on changes in the same section.}
Two, both mechanical.
\begin{itemize}\itemsep2pt
    \item The definition of $i^L_K$ is currently phrased through a uniformizer
    $\Pi$ with $B=A[\Pi]$. Under the weaker hypothesis the monogenic generator
    need not be a uniformizer, so it becomes
    \new{$i^L_K(\sigma)=v_B(\sigma(x)-x)$} for the element $x$ with $B=A[x]$.
    Nothing downstream depends on $x$ being a uniformizer: what is used is that
    $G_i=\{\sigma: i^L_K(\sigma)\ge i+1\}$.
    \item The parenthetical ``(complete local fields with a discrete valuation
    and perfect residue field)'' becomes \new{``\dots and separable residue
    extension''}.
\end{itemize}

\paragraph{Safe downstream?}
Yes. The section on the cyclic theories re-declares its own hypothesis
(``let $K$ be a complete discrete valuation field with perfect residue field
$\kappa$ of characteristic $p>0$''), so the Artin--Schreier--Witt and Kummer
computations keep the explicit $k((t))$ description they need.

% =====================================================================
\chapter*{Edit B --- the general definition}
\addcontentsline{toc}{chapter}{Edit B --- the definition}
% =====================================================================
\loc{File \texttt{sections/Preliminaries/GTorsorRamification.tex}, inserted after
\Cref{definition:tamely_ramified_in_codim_1} ends (line 24) --- that is,
\emph{before} ``Now back to our original assumptions''.}

The point of the placement: the section already sets up $X$, $U$, $(Z,\xi)$ in
the right generality and then abandons it. This edit finishes the general setup
before descending to the curve, rather than adding a new subsection.

\begin{proposed}
\new{\Cref{definition:tamely_ramified_in_codim_1} is qualitative. To measure
ramification we specialize to torsors and pass to the completion. Keep $X$, $U$
and $(Z,\xi)$ as above, let $G$ be a finite abelian group, and let
$\cP\to U$ be a $G$-torsor in $U_{\text{\'Et}}$. Write $\widehat{\Oo}_{X,\xi}$
for the completion of $\Oo_{X,\xi}$ and
$\widehat{K}_\xi=\Frac(\widehat{\Oo}_{X,\xi})$ for the associated complete
discretely valued field. Restricting $\cP$ to the punctured formal
neighbourhood $\Spec(\widehat{K}_\xi)$ isolates the ramification at $(Z,\xi)$:
this restriction is finite \'etale over $\widehat{K}_\xi$, hence by
\stackstag{02GL} decomposes into a finite disjoint union of spectra of finite
separable field extensions,}
\[
    \new{\cP|_{\Spec(\widehat{K}_\xi)}\cong\bigsqcup_{i=1}^{N}\Spec(F_i),}
\]
\new{with the following properties:}
\begin{enumerate}\color{newmath}
    \item The extensions $F_i/\widehat{K}_\xi$ are mutually isomorphic and
    Galois, and each $F_i$ is complete for the unique extension of the discrete
    valuation on $\widehat{K}_\xi$
    \cite[Chapter~II, \S 2, Proposition~3]{serreLF}.
    \item If $F$ denotes any one of the $F_i$, then
    $H=\Gal(F/\widehat{K}_\xi)$ identifies with the stabilizer of the
    corresponding connected component of $\cP|_{\Spec(\widehat{K}_\xi)}$, which
    is a subgroup of $G$.
    \item The number of components is $N=[G:H]$.
\end{enumerate}

\begin{definition}\label{definition:local_ramification_boundness}
\color{newmath}
    In the situation above, assume that the residue extension of
    $F/\widehat{K}_\xi$ is separable, so that the upper-numbering filtration of
    \Cref{section:ramification_of_discrete_valuations} is defined for
    $F/\widehat{K}_\xi$; we then call $(Z,\xi)$ \emph{residually separable for
    $\cP$}. We say that the extension $F/\widehat{K}_\xi$ \textit{has
    ramification bounded by $r$} if the ramification group $H^r$ is trivial.
    Here $H^r$ denotes the upper-numbering ramification group of
    $H=\Gal(F/\widehat{K}_\xi)$ associated with the extension
    $F/\widehat{K}_\xi$; it is not a filtration inherited from the ambient
    group $G$. This condition is independent of the chosen connected component
    of $\cP|_{\Spec(\widehat{K}_\xi)}$, and we say \textit{$\cP$ has
    ramification bounded by $r$ at $(Z,\xi)$} when it holds.
\end{definition}

\new{It is convenient to restate this through characters. Since $G$ is a finite
abelian group over $k$, hence \'etale, the correspondence between $G$-torsors
and the fundamental group
(\Cref{cor:abelian_equivilance_fundamental_torsors}) applies; over the complete
field $\widehat{K}_\xi$ the fundamental group is the absolute Galois group
$G_{\widehat{K}_\xi}:=\Gal(\widehat{K}_\xi^{\mathrm{sep}}/\widehat{K}_\xi)$, so
that $\cP|_{\Spec(\widehat{K}_\xi)}$ corresponds to a unique continuous
homomorphism}
\[
    \new{\rho_\xi:G_{\widehat{K}_\xi}\longrightarrow G,}
\]
\new{with $H=\operatorname{im}(\rho_\xi)$ and
$F=(\widehat{K}_\xi^{\mathrm{sep}})^{\ker\rho_\xi}$. Because the upper
numbering is compatible with quotients, $H^r$ is the image of
$G^{\,r}_{\widehat{K}_\xi}$, and therefore}
\begin{equation}\label{eq:ramification_bound_via_character}
    \new{\cP\ \text{has ramification bounded by } r \text{ at } (Z,\xi)
    \quad\Longleftrightarrow\quad
    \rho_\xi\bigl(G_{\widehat{K}_\xi}^{\,r}\bigr)=\{1\}.}
\end{equation}

\begin{lemma}\label{lemma:bound_zero_one_versus_codim_one}
\color{newmath}
    Let $(Z,\xi)$ be residually separable for $\cP$. Then $\cP$ has ramification
    bounded by $0$ at $(Z,\xi)$ if and only if it is unramified there, and
    bounded by $1$ if and only if it is tamely ramified there, both in the sense
    of \Cref{definition:tamely_ramified_in_codim_1}.
\end{lemma}
\begin{proof}
\color{newmath}
    Ramification indices and residue extensions are unchanged under the
    faithfully flat extension $\Oo_{X,\xi}\to\widehat{\Oo}_{X,\xi}$, which
    induces an isomorphism on residue fields; so both conditions of
    \Cref{definition:tamely_ramified_in_codim_1} may be tested after
    completion. Now $H^0=H_0$ is the inertia group, which gives the first
    claim. For the second, $H_1=1$ implies $H^1=1$ directly. Conversely,
    since $H\subseteq G$ is abelian, the jumps of the upper-numbering
    filtration are integers by the Hasse--Arf theorem
    \cite[Chapter~V, \S 7, Theorem~1]{serreLF}; the jump at which $H_1$ dies
    therefore occurs at an integral upper index, so $H^1=1$ forces $H_1=1$.
\end{proof}
\end{proposed}

\paragraph{Why the character paragraph is the load-bearing half.}
\Cref{eq:ramification_bound_via_character} is the form in which every proof in
the subsection actually uses the notion. With it stated at a general $\xi$, the
existing proofs of the contracted-product lemma, the \'etale-descent lemma, the
fibre-product lemma and the induced-torsor lemma go through
\emph{without a single word changed} --- they are already written in this
language.

% =====================================================================
\chapter*{Edit C --- collapse the curve paragraph}
\addcontentsline{toc}{chapter}{Edit C --- the curve case}
% =====================================================================
\loc{Same file, replacing lines 26--57 (the closed-point decomposition and the
old \emph{ramification bounded by $r$} definition), and deleting lines
155--173.}

\begin{current}
Now back to our original assumptions, $C$ is a projective smooth geometrically
connected curve over a field $k$ and $U=C\setminus\mdls$. Every $(Z,\xi)$ is now
a closed point $P$ \dots\ [followed by the decomposition
$\cP|_{\Spec(\widehat K_P)}\cong\bigsqcup\Spec(F_i)$, its three properties, and
the definition of \emph{ramification bounded by $r$ at $P$}].
\end{current}

\begin{proposed}
Now back to our original assumptions: $C$ is a projective smooth geometrically
connected curve over a field $k$ and $U=C\setminus\mdls$. Every $(Z,\xi)$ is now
a closed point $P$, the local ring $\Oo_{C,P}$ is a discrete valuation ring, and
$\widehat{K}_P=\Frac(\widehat{\Oo}_{C,P})$ is the complete local field at $P$.
\new{The residual separability hypothesis of
\Cref{definition:local_ramification_boundness} is automatic here, since
$\kappa(P)$ is finite over the perfect field $k$ and hence perfect. We write
\textit{$\cP$ has ramification bounded by $r$ at $P$} for the condition of
\Cref{definition:local_ramification_boundness} applied to
$(Z,\xi)=(\{P\},P)$.}
\end{proposed}

The two modulus definitions and the proposition comparing the closed-point and
geometric formulations then follow \emph{unchanged}.

\paragraph{Two deletions.}
\begin{itemize}\itemsep2pt
    \item The sentence ``Furthermore, the notions of tame ramification and
    unramifiedness derived here coincide with the codimension-$1$ criteria in
    \Cref{definition:tamely_ramified_in_codim_1}'' is now
    \Cref{lemma:bound_zero_one_versus_codim_one}, with a proof.
    \item The paragraph beginning ``There is another equivalent formulation
    through the language of characters'' moves upstream into Edit~B and is
    deleted here.
\end{itemize}

% =====================================================================
\chapter*{Edit D --- pin the lemma hypotheses}
\addcontentsline{toc}{chapter}{Edit D --- lemma hypotheses}
% =====================================================================
\loc{Same file, the statements of the four stability lemmas.}

The contracted-product lemma currently opens with a free-standing setup. It
should point at the definition, exactly as the fibre-product lemma already
points at it:

\begin{proposed}
Let $G$ be a finite abelian group, and let $X$, $U$ and $(Z,\xi)$ be as in
\new{\Cref{definition:local_ramification_boundness}, with $(Z,\xi)$ residually
separable for both $\cP_1$ and $\cP_2$.} Assume $\cP_1$ has ramification
bounded by $r_1$ at $(Z,\xi)$ and $\cP_2$ has ramification bounded by $r_2$ at
$(Z,\xi)$. Then their contracted product $\cP_1\otimes^G\cP_2$ has ramification
bounded by $\max(r_1,r_2)$ at $(Z,\xi)$.
\end{proposed}

The same clause is added to the \'etale-descent lemma, at both $\xi_X$ and
$\xi_Y$ --- with the remark that there residual separability at $\xi_X$
\emph{follows} from residual separability at $\xi_Y$, since $f$ is \'etale at
$\xi_X$, so the hypothesis need only be imposed downstairs. The fibre-product
and induced-torsor lemmas already say ``let $X,U,(Z,\xi)$ be as in'' the
contracted-product lemma and so inherit it automatically.

% =====================================================================
\chapter*{Consequences, and three things to decide}
\addcontentsline{toc}{chapter}{Consequences and open points}
% =====================================================================

\paragraph{Numbering.}
Adding \Cref{lemma:bound_zero_one_versus_codim_one} as a numbered statement
shifts everything after it by one: the contracted-product lemma moves from
Lemma~43 to Lemma~44, and so on to the end of the thesis. Every reference is
\verb|\Cref|'d, so this recompiles correctly, but ``Lemma 43'' stops meaning
what it means today. The alternative is to set it as a \verb|rem*| --- your
\texttt{remarkstyle} is \texttt{numbered=no} --- which preserves all numbering
at the price of not being referenceable.

\paragraph{Where residual separability must now be checked.}
The hypothesis is genuinely needed at $\eta_\ndls$, and it is genuinely
verified there: the Artin--Schreier argument shows each factor of
$B\otimes_A\widehat K_\eta$ has ``ramification index $1$ and separable residue
extension''. What the edit changes is that this verification now discharges a
stated hypothesis instead of being an unmotivated step. Worth a one-line
forward reference from the definition to that computation.

\paragraph{One reference to fill in.}
The claim that ramification indices and residue extensions are unchanged by
completion is argued in one sentence in the proof of
\Cref{lemma:bound_zero_one_versus_codim_one}. It is elementary
($B\otimes_A\widehat A\cong\prod_i\widehat{B_{\mathfrak m_i}}$, preserving $e$
and $f$), but if you would prefer a citation there I would rather look up the
Stacks tag than guess at one.

\paragraph{Alternative worth a moment's thought.}
Everything above keeps ``bounded by $r$'' anchored to the upper-numbering
filtration of the finite quotient $H$. One could instead take
\eqref{eq:ramification_bound_via_character} --- the condition
$\rho_\xi(G^{\,r}_{\widehat K_\xi})=1$ on the absolute Galois group --- as the
\emph{definition}, and derive the $H^r$ description as a consequence. That
reads more cleanly, matches Guignard's formulation, and makes the stability
lemmas nearly tautological; the cost is that it obscures the elementary content
and makes the comparison with
\Cref{definition:tamely_ramified_in_codim_1} less direct. I did not take it,
but it is a real choice and it is easy to switch.

\begin{thebibliography}{9}
\bibitem{serreLF} J.-P.~Serre, \emph{Corps Locaux} (\emph{Local Fields}),
Hermann / Springer GTM~67.
\bibitem{stacks-project} The Stacks Project Authors, \emph{The Stacks Project},
\url{https://stacks.math.columbia.edu}.
\bibitem{Guignard2018} Q.~Guignard, \emph{Geometric local class field theory}.
\end{thebibliography}

\end{document}